% ---------------------------------------------------------------------------
% Ameliorating inventory section.
%
% INPUT-ABLE FRAGMENT: this file is \input{} into the manuscript
% learning_inventory_control_policies_es.tex and reuses ITS preamble macros and
% TikZ styles (envbox, envtitle, envsupply/envstock/envsink/envlost, reqflow/
% matflow/demandflow/lostflow, ltlabel/reqlabel, costH/costP/costK, the boxed
% "Per-period event sequence" panel, booktabs, \ding). It defines NO new styles
% and contains NO \documentclass.
%
% BIBLIOGRAPHY: this section cites \citet{pahrGrunow2025}, which is NOT yet in
% paper/references.bib. The BibTeX entry is provided inline in
% scripts/ameliorating_inventory/PAPER_SECTION_DRAFT.md and must be added to
% references.bib before compiling (out of scope of this fragment).
%
% ALGORITHMIC / DATA-PROVENANCE NOTE (every number is sourced from the repo):
%   Instance + dynamics : src/problems/ameliorating_inventory/
%       average_profit_blending_env.rs (step_state: purchase -> age with
%       evaporation + age-dependent beta decay -> per-period blending LP solves
%       issuance, production derived -> sales; long-run average profit),
%       perfect_information_lp.rs (re-solves the perfect-information LP; executing
%       test reproduces both published upper bounds to <1e-3 in tests/
%       verification.rs), references.rs (spirits_0001, port_wine configs).
%   Comparator          : (i) the perfect-information LP UPPER BOUND on average
%       profit (1991.9344 spirits_0001, 2444.8011 port_wine) -- a bound, report
%       gap, NEVER "beat"; (ii) the in-repo tuned order-up-to purchase gate --
%       beatable.
%   Results             : outputs/autoresearch/ameliorating_inventory_average_
%       profit_autoresearch/{spirits_0001,port_wine}_d1_oblique_full.json +
%       scripts/ameliorating_inventory/RESULTS_FULL_BUDGET.md (committed 67e44e4,
%       full budget). spirits_0001 learned 115.07 +/- 0.44 vs gate 20.91
%       (+94.16, +450%, ~140x paired SEM; gap-to-bound 94.2%); port_wine learned
%       505.78 +/- 0.59 vs gate 133.78 (+372.00, +278%, ~610x paired SEM;
%       gap-to-bound 79.3%).
%   HONEST FRAMING      : learned BEATS the order-up-to gate; the gap to the LP
%       bound is STRUCTURAL (single scalar purchase action vs the bound's full
%       3-part LP issuance + ~200/unit purchase cost every period) and is NOT
%       comparable to the paper's ~3.5% DRL gap (full 3-part action).
% ---------------------------------------------------------------------------
\section{Ameliorating inventory}
\label{sec:amelior}

\subsection{Problem formulation}
\label{sec:amelior-formulation}

Figure~\ref{fig:env-amelior} depicts this environment as a flow network --- the price-dependent
purchase into the youngest age class, the aging of stock that \emph{improves} with age (subject to
evaporation and stochastic decay), the per-period blending and issuance to age-targeted products, and
where purchase cost and holding are charged against sales revenue.

\begin{figure}[t]
  \centering
  \resizebox{\textwidth}{!}{%
  \begin{tikzpicture}[x=1cm,y=1cm]
    % Single-item ameliorating system: product IMPROVES with age. Purchase at a
    % stochastic price enters the youngest age class; stock ages (evaporation +
    % age-dependent decay to a loss sink); a per-period blending LP issues aged
    % stock to products defined by TARGET AGES, sold at stochastic prices.
    \node[envsupply] (sup)  at (0,0)      {external\\supplier\\(price $P_t$)};
    \node[envstock]  (stk)  at (7.4,0)    {age-structured stock\\$\mathbf I_t=(I^{1}_t,\dots,I^{A}_t)$\\(young $\to$ aged/improved)};
    \node[envsink]   (prod) at (15.6,0)   {age-targeted\\products (sales $P^{s}_t$)};
    \node[envlost]   (loss) at (7.4,-4.2) {evaporation\\$+$ decay\\\emph{(loss)}};

    % purchase: dashed request UP, solid material DOWN into youngest class
    \draw[reqflow] (stk.north west) to[bend right=32] (sup.north east);
    \node[reqlabel] at (3.3,2.0) {\ding{183}\, purchase $a^P_t$};
    \draw[matflow] (sup.south east) to[bend right=10] (stk.west);
    \node[ltlabel]  at (3.4,-1.55) {\ding{184}\, enters youngest class;\\cost $P_t\,a^P_t$};

    % issue/blend to products (demand-style), decay to loss sink
    \draw[demandflow] (stk) -- (prod)
          node[midway,above=4pt,ltlabel] {\ding{186}\, blending LP issues\\aged stock to products\\(mean age $\ge$ target)};
    \draw[lostflow,myorange!70!black] (stk.south) -- (loss.north);

    \node[costH,above=9mm of stk.north] {holding $h\sum_a I^{a}_t$};
    \node[costK,below=2pt of loss.south,align=center] {ameliorate:\\age $\Rightarrow$ value};
    \node[costP,right=3pt of prod.east] {revenue $+\,P^{s}_t$};

    \begin{scope}[on background layer]
      \node[envbox,fit=(current bounding box)] (box) {};
    \end{scope}
    \node[envtitle,above=3pt of box.north,anchor=south]
          {Ameliorating single-item environment ($A$ age classes; long-run average profit)};
  \end{tikzpicture}%
  }
  \par\vspace{6pt}\centerline{%
    \tikz\node[draw=mydarkblue!55,rounded corners=4pt,fill=mydarkblue!4,inner sep=7pt]{%
      \begin{minipage}{\dimexpr\linewidth-18pt\relax}\centering
        {\small\bfseries\color{mydarkblue!85}Per-period event sequence}\\[4pt]
        {\small\begin{tabular}{@{}p{4.3cm}p{4.3cm}p{4.3cm}@{}}
          \ding{182}\,observe inventory by age $\mathbf I_t$ and realized purchase price $P_t$ & \ding{183}\,choose purchase volume $a^P_t\in[0,\bar I]$; it enters the youngest age class & \ding{184}\,a per-period \emph{blending LP} issues aged stock to products whose blend mean age meets the target\\[4pt]
          \ding{185}\,sales realize at stochastic prices $P^{s}_t$; production is derived from the LP & \ding{186}\,stock \emph{ages}: evaporation and age-dependent decay remove a fraction (loss) & \ding{187}\,profit $=$ sales revenue $-$ purchase cost $P_t a^P_t-$ holding $h\sum_a I^{a}_t$\\
        \end{tabular}}
      \end{minipage}};}\par
  \caption{Ameliorating-inventory environment of \citet{pahrGrunow2025}, drawn as an event-based
  interaction ($A$ age classes; the product \emph{improves} with age, e.g.\ spirits or port wine).
  Each period the controller \emph{purchases} $a^P_t$ units at a realized stochastic price $P_t$ (thin
  dashed control arrow); the goods enter the \emph{youngest} age class (solid material arrow, cost
  $P_t a^P_t$). A per-period \emph{blending LP} then issues aged stock to products defined by
  \emph{target ages} --- a product can ship only if the issued blend's mean age meets its target ---
  and sales realize at stochastic prices $P^{s}_t$; production is derived from the LP solution. Stock
  \emph{ages} subject to evaporation and an age-dependent stochastic decay (orange loss sink). The
  objective is long-run \emph{average profit}: sales revenue minus purchase cost and holding; the
  boxed panel gives the per-period event sequence (Section~\ref{sec:amelior-formulation},
  Eq.~\eqref{eq:amelior-profit}).}
  \label{fig:env-amelior}
\end{figure}

We consider the single-item, periodic-review ameliorating-inventory control problem of
\citet{pahrGrunow2025}, in which the product \emph{improves} with age: inventory is purchased young
and ages through $A$ discrete age classes, gaining value (and meeting product target ages) but subject
to evaporation and stochastic decay. The objective is long-run \emph{average profit}. We use the
faithful average-profit formulation, including the per-period blending linear program (LP) that solves
issuance and from which production is derived.

\paragraph{Timing of a period.}
At the start of period $t$ the controller observes the inventory disaggregated by age class and the
realized purchase price $P_t$, and chooses a purchase volume $a^P_t$ which enters the youngest age
class. A per-period blending LP issues aged stock to the products --- each product is defined by a
target age and (with blending) may be served by a mix of age classes whose mean age meets the target
--- and sales realize at stochastic prices. The remaining stock then ages: a fraction is lost to
evaporation and an age-dependent stochastic decay.

\paragraph{State.}
The state is the inventory vector disaggregated by age class together with the realized purchase
price,
\begin{equation*}
    \mathbf{S}_t=\bigl(I^{1}_t,\dots,I^{A}_t;\ P_t\bigr),
\end{equation*}
where $I^{a}_t$ is the on-hand quantity in age class $a$ ($a=1$ youngest, $a=A$ oldest/most improved).
The sales prices follow a correlated process and form part of the environment's stochasticity.

\paragraph{Action.}
In the faithful dynamics the only free control is the scalar \emph{purchase} volume
\begin{equation*}
    a_t=a^P_t\in[0,\bar I],
\end{equation*}
where $\bar I$ is the inventory capacity. Issuance is determined by the environment's per-period
blending LP and production is derived from it, so the ``three-part'' purchase/produce/issue decision
of \citet{pahrGrunow2025} is reduced here to the structural purchase head.

\paragraph{Transition.}
The purchase enters the youngest class, $I^{1}\!\leftarrow a^P_t$. The blending LP issues quantities
$x^{a}_t$ from the age classes to meet product demand subject to the per-product mean-age targets and a
processing capacity; sold quantities leave inventory. The surviving stock then ages by one class, and
an age-dependent decay factor together with a constant evaporation removes a fraction of each class
(the loss sink in Figure~\ref{fig:env-amelior}).

\paragraph{One-period profit.}
Writing $R_t$ for realized sales revenue (issued quantities valued at the stochastic sales prices),
the profit incurred in period $t$ is revenue minus the purchase cost and holding:
\begin{equation}
    \label{eq:amelior-profit}
    \rho_t \;=\; R_t \;-\; P_t\,a^P_t \;-\; h\sum_{a=1}^{A} I^{a}_t ,
\end{equation}
with stochastic purchase price $P_t$ and per-unit holding cost $h$.

\paragraph{Objective.}
A stationary policy $\pi_{\btheta}$ maps the state to the purchase volume. Performance is the long-run
\emph{average profit},
\begin{equation}
    \label{eq:amelior-objective}
    G(\btheta) \;=\; \lim_{T\to\infty}\frac{1}{T}\,\mathbb{E}\!\left[\sum_{t=1}^{T}\rho_t\right],
    \qquad \max_{\btheta}\ G(\btheta),
\end{equation}
estimated by simulation over shared price/demand-path seeds on a long post-warm-up window.

\FloatBarrier

\subsection{Policy}
\label{sec:amelior-policy}

Ameliorating inventory tests the recipe under a \emph{profit} objective with a price-augmented state.
\textbf{Input.} The policy input is the price-augmented state $[\,P_t,\ I^{1}_t,\dots,I^{A}_t\,]$,
divide-by-scale normalized by the inventory capacity. \textbf{Output.} A valid action is the scalar
purchase volume $a^P_t\in[0,\bar I]$. \textbf{Internals.} $\pi_{\btheta}$ uses the soft-tree backbone
of Section~\ref{sec:soft-tree} with a single purchase head and a \emph{linear} leaf, so the head can
express a \emph{price-reactive} order-up-to purchase --- buying more when the realized purchase price
is low, which a fixed order-up-to level cannot do. We warm-start CMA-ES at an order-up-to purchase
$a^P_t=\mathrm{softplus}\!\left(S-\sum_a I^{a}_t\right)$, so generation~0 reproduces a simple
order-up-to heuristic and the search refines a price-reactive purchase.

\subsection{Experiment setup}
\label{sec:amelior-setup}

We use the two reference instances of Table~\ref{tab:amelior-params}: the companion default
\texttt{spirits\_0001} (10 age classes, 3 products, no blending) and the industry \texttt{port\_wine}
case study (25 age classes, 2 products, blending enabled).

\paragraph{Comparators: a bound and a beatable gate.}
This problem has two reference points of different kinds, and we treat them differently. (i)~The
\emph{perfect-information LP upper bound} on average profit ($1991.93$ for \texttt{spirits\_0001},
$2444.80$ for \texttt{port\_wine}) is an \emph{upper bound} under hindsight and full LP issuance: our
implementation \emph{re-solves} this LP and reproduces both published bounds to within $10^{-3}$ (an
executing fidelity test, not a stored constant), and we report the \emph{gap} to it but never treat it
as an achievable target. (ii)~The like-for-like keep/discard \emph{gate} is the best tuned order-up-to
purchase (the level $S$ maximizing profit on the same evaluation block); beating it is a legitimate
win.

\paragraph{Evaluation protocol.}
Policies are trained with warm-started CMA-ES (full budget: population $16$, $60$ generations) and
evaluated under the common-random-number protocol of Section~\ref{sec:methods-eval} on a held-out
block ($12{,}000$-period horizon, $24$ paired seeds disjoint from training). A win over the gate is
claimed only when the paired-difference advantage exceeds its standard error.

\begin{table}[!htbp]
\centering
\small
\caption{Ameliorating-inventory reference instances of \citet{pahrGrunow2025} used here. Source:
\texttt{src/problems/ameliorating\_inventory/references.rs} and the checked-in perfect-information LP
datasets.}
\label{tab:amelior-params}
\begin{tabular}{llll}
\toprule
Instance & Age classes $A$ & Products (target ages) & Perfect-info LP bound \\
\midrule
\texttt{spirits\_0001} & $10$ & $3$ (ages $\{2,4,6\}$), no blending & $1991.93$ \\
\texttt{port\_wine}    & $25$ & $2$ (ages $\{9,19\}$), blending     & $2444.80$ \\
\bottomrule
\end{tabular}
\end{table}
\FloatBarrier

\subsection{Results}
\label{sec:amelior-results}

Table~\ref{tab:amelior-results} reports the learned price-reactive purchase policy against the best
tuned order-up-to gate at full budget. On both instances the learned policy \emph{beats the gate by a
wide, statistically unambiguous margin}: on \texttt{spirits\_0001} it raises average profit from
$20.91$ to $115.07$ ($+94.16$, $+450\%$, $\approx140\times$ the paired SEM), and on \texttt{port\_wine}
from $133.78$ to $505.78$ ($+372.00$, $+278\%$, $\approx610\times$ the paired SEM). The mechanism is
price reactivity: the linear leaf buys more when the realized purchase price is low, which a fixed
order-up-to level cannot exploit.

\paragraph{Interpretation: an honest gap to an upper bound.}
The gap to the perfect-information LP bound remains large ($94.2\%$ on \texttt{spirits\_0001},
$79.3\%$ on \texttt{port\_wine}) and we report it as such. This gap is \emph{structural}, not an
optimization failure: the bound assumes perfect information and the full three-part LP decision
(purchase, production targets, and per-age issuance solved jointly with hindsight), whereas our policy
controls only the scalar purchase volume while the environment charges the full purchase price
($\sim$200 per unit) every period. A feasible single-purchase policy on the stochastic environment
therefore sits well below the hindsight bound by construction, and our gap is \emph{not} comparable to
the $\sim$3.5\% gap reported by \citet{pahrGrunow2025} for a deep-RL agent that acts in the full
three-part action space. The tighter gap on \texttt{port\_wine} ($79.3\%$ vs $94.2\%$) is explained by
blending: issuing across the two target ages lets each purchased unit convert into more sold product,
so a single purchase lever captures more of the bound. The honest claim is therefore a robust beat of
the in-repo order-up-to heuristic under a single shared estimator; widening the action geometry to the
full three-part decision is the natural next step toward the deep-RL gap and is left to future work.

\begin{table}[!htbp]
\centering
\small
\caption{Ameliorating inventory: best learned price-reactive purchase soft tree vs.\ the best tuned
order-up-to gate, common-random-number evaluation ($12{,}000$-period horizon, $24$ paired seeds).
Higher average profit is better. ``Gain over gate'' is reported as the absolute profit gain and as a
fraction of the gate; ``gap to bound'' is the fraction of the perfect-information LP \emph{upper bound}
left uncaptured (the bound is a bound, never an achievable target). Source:
\texttt{outputs/autoresearch/ameliorating\_inventory\_average\_profit\_autoresearch/}.}
\label{tab:amelior-results}
\begin{tabular}{lccccc}
\toprule
Instance & Learned ($\pm$ SEM) & Order-up-to gate & Gain over gate & LP bound & Gap to bound \\
\midrule
\texttt{spirits\_0001} & $\mathbf{115.07\pm0.44}$ & $20.91$  & $+94.16\ (+450\%)$ & $1991.93$ & $94.2\%$ \\
\texttt{port\_wine}    & $\mathbf{505.78\pm0.59}$ & $133.78$ & $+372.00\ (+278\%)$ & $2444.80$ & $79.3\%$ \\
\bottomrule
\end{tabular}
\par\medskip\footnotesize The like-for-like win is over the tuned order-up-to gate (paired CRN,
$\approx140\times$ / $\approx610\times$ SEM). The LP-bound gap is structural --- a single scalar
purchase action vs the bound's full three-part LP issuance --- and is not comparable to the
$\sim$3.5\% deep-RL gap of \citet{pahrGrunow2025}; it is reported, never ``beaten''.
\end{table}
\FloatBarrier
